% GENERATED FILE. Edit canonical lessons, then run npm run generate:books.
% canonical-source-hash: fnv1a64:00d5619d5abc84d9
% canonical-lessons: HI-C521-ac, HI-C521-dalna, HI-C521-jira, HI-C521-bagharna, HI-C521-bhujna, HI-R521-review-514-521

\chapter{Heat in the Kitchen}
\label{ch:hi-heat-in-the-kitchen}
\hlchaptermodality{521}

\begin{chapteropening}
\textbf{By the end of this chapter:} \emph{I can talk about the flame, sizzling spices and steam.}

Complete the last lesson of chapter 521: I can talk about the flame, sizzling spices and steam.
\end{chapteropening}

\section[\texorpdfstring{\hi{आँच} (ā̃c)}{आँच (ā̃c)}]{\hi{आँच} (ā̃c) — the heat of a flame; a flame}

\label{lesson:HI-C521-ac}



\begin{quote}
Before the new one: say the Hindi for to roll out (dough), then the Hindi for to puff up.
\end{quote}

\subsection*{You'll want to know: \hi{आँच}}
\textbf{\hi{आँच}} — \emph{ā̃c} — \textquotedblleft{}the heat of a flame; a flame\textquotedblright{}.

\begin{quote}
\hi{दाल} \hi{को} \hi{धीमी} \hi{आँच} \hi{पर} \hi{पकाइए।}
\end{quote}

\begin{quote}
\emph{Dāl ko dhīmī ā̃c par pakāie.}
\end{quote}

\begin{quote}
Cook the dal on a low flame.
\end{quote}

Feminine.

\begin{cousinweb}
\textbf{\hi{आँच}} (\emph{ā̃c}) is a tadbhava, inherited from Sanskrit \emph{arcis}, \textquotedblleft{}a flame, a ray\textquotedblright{} (CDIAL 635).
\end{cousinweb}

\begin{grammarlens}[title={What you've built}]
One more word for describing what you lived and felt.
\end{grammarlens}

\subsection*{Your turn}
\begin{itemize}
\raggedright
  \item \emph{Say it:} \emph{ā̃c}
  \item \emph{Say it:} \emph{ā̃c}, once more
  \item \emph{Say it:} \emph{Dāl ko dhīmī ā̃c par pakāie.} — Cook the dal on a low flame.
  \item \emph{Recall:} say the Hindi for to roll out (dough), then the Hindi for to puff up, then say \emph{ā̃c} again
\end{itemize}

\begin{tcolorbox}[breakable,colback=teal!4,colframe=teal!35!black,title={Before you move on}]
What does \hi{आँच} mean? (\textquotedblleft{}The heat of a flame; a flame\textquotedblright{}.) Say it once more.
\end{tcolorbox}
\section[\texorpdfstring{\hi{डालना} (ḍālnā)}{डालना (ḍālnā)}]{\hi{डालना} (ḍālnā) — to put in; to pour in}

\label{lesson:HI-C521-dalna}



\begin{quote}
Before the new one: say the Hindi for to puff up, then the Hindi for the heat of a flame.
\end{quote}

\subsection*{You'll want to know: \hi{डालना}}
\textbf{\hi{डालना}} — \emph{ḍālnā} — \textquotedblleft{}to put in; to pour in\textquotedblright{}.

\begin{quote}
\hi{गरम} \hi{तेल} \hi{में} \hi{थोड़ा} \hi{नमक} \hi{डालिए।}
\end{quote}

\begin{quote}
\emph{Garam tel meṁ thoṛā namak ḍāliye.}
\end{quote}

\begin{quote}
Put a little salt into the hot oil.
\end{quote}

Past, with \hi{ने} (he or she …): \hi{उसने} \hi{दाल} \hi{में} \hi{पानी} \hi{डाला}, \emph{usne dāl meṁ pānī ḍālā}.

\begin{cousinweb}
\textbf{\hi{डालना}} (\emph{ḍālnā}) is a tadbhava, inherited from \emph{*ḍāl}, \textquotedblleft{}pour, put\textquotedblright{}, a form CDIAL reconstructs from its descendants (CDIAL 5545).
\end{cousinweb}

\begin{grammarlens}[title={What you've built}]
One more word for describing what you lived and felt.
\end{grammarlens}

\subsection*{Your turn}
\begin{itemize}
\raggedright
  \item \emph{Say it:} \emph{ḍālnā}
  \item \emph{Say it:} \emph{ḍālnā}, once more
  \item \emph{Say it:} \emph{Garam tel meṁ thoṛā namak ḍāliye.} — Put a little salt into the hot oil.
  \item \emph{Recall:} say the Hindi for to puff up, then the Hindi for the heat of a flame, then say \emph{ḍālnā} again
\end{itemize}

\begin{tcolorbox}[breakable,colback=teal!4,colframe=teal!35!black,title={Before you move on}]
What does \hi{डालना} mean? (\textquotedblleft{}To put in; to pour in\textquotedblright{}.) Say it once more.
\end{tcolorbox}
\section[\texorpdfstring{\hi{जीरा} (jīrā)}{जीरा (jīrā)}]{\hi{जीरा} (jīrā) — cumin}

\label{lesson:HI-C521-jira}



\begin{quote}
Before the new one: say the Hindi for the heat of a flame, then the Hindi for to put in.
\end{quote}

\subsection*{You'll want to know: \hi{जीरा}}
\textbf{\hi{जीरा}} — \emph{jīrā} — \textquotedblleft{}cumin\textquotedblright{}.

\begin{quote}
\hi{माँ} \hi{ने} \hi{गरम} \hi{घी} \hi{में} \hi{जीरा} \hi{डाला।}
\end{quote}

\begin{quote}
\emph{Māṁ ne garam ghī meṁ jīrā ḍālā.}
\end{quote}

\begin{quote}
Mother put cumin into the hot ghee.
\end{quote}

Masculine.

\begin{cousinweb}
\textbf{\hi{जीरा}} (\emph{jīrā}) is a tadbhava, inherited from Sanskrit \emph{jīraka}, \textquotedblleft{}cumin\textquotedblright{} (CDIAL 5234).
\end{cousinweb}

\begin{grammarlens}[title={What you've built}]
One more word for describing what you lived and felt.
\end{grammarlens}

\subsection*{Your turn}
\begin{itemize}
\raggedright
  \item \emph{Say it:} \emph{jīrā}
  \item \emph{Say it:} \emph{jīrā}, once more
  \item \emph{Say it:} \emph{Māṁ ne garam ghī meṁ jīrā ḍālā.} — Mother put cumin into the hot ghee.
  \item \emph{Recall:} say the Hindi for the heat of a flame, then the Hindi for to put in, then say \emph{jīrā} again
\end{itemize}

\begin{tcolorbox}[breakable,colback=teal!4,colframe=teal!35!black,title={Before you move on}]
What does \hi{जीरा} mean? (\textquotedblleft{}Cumin\textquotedblright{}.) Say it once more.
\end{tcolorbox}
\section[\texorpdfstring{\hi{बघारना} (baghārnā)}{बघारना (baghārnā)}]{\hi{बघारना} (baghārnā) — to temper (with hot spiced oil)}

\label{lesson:HI-C521-bagharna}



\begin{quote}
Before the new one: say the Hindi for to put in, then the Hindi for cumin.
\end{quote}

\subsection*{You'll want to know: \hi{बघारना}}
\textbf{\hi{बघारना}} — \emph{baghārnā} — \textquotedblleft{}to temper (with hot spiced oil)\textquotedblright{}.

\begin{quote}
\hi{माँ} \hi{ने} \hi{दाल} \hi{को} \hi{घी} \hi{और} \hi{जीरे} \hi{से} \hi{बघारा।}
\end{quote}

\begin{quote}
\emph{Māṁ ne dāl ko ghī aur jīre se baghārā.}
\end{quote}

\begin{quote}
Mother tempered the dal with ghee and cumin.
\end{quote}

Past, with \hi{ने} (he or she …): \hi{उसने} \hi{सब्ज़ी} \hi{बघारी}, \emph{usne sabzī baghārī}.

\begin{cousinweb}
\textbf{\hi{बघारना}} (\emph{baghārnā}) is a tadbhava, inherited from Sanskrit \emph{vyāghārayati}, \textquotedblleft{}sprinkles (with hot fat)\textquotedblright{} (CDIAL 12191). CDIAL gives the Hindi verb as \textquotedblleft{}to brown onions and spices in heated oil\textquotedblright{}.
\end{cousinweb}

\begin{grammarlens}[title={What you've built}]
One more word for describing what you lived and felt.
\end{grammarlens}

\subsection*{Your turn}
\begin{itemize}
\raggedright
  \item \emph{Say it:} \emph{baghārnā}
  \item \emph{Say it:} \emph{baghārnā}, once more
  \item \emph{Say it:} \emph{Māṁ ne dāl ko ghī aur jīre se baghārā.} — Mother tempered the dal with ghee and cumin.
  \item \emph{Recall:} say the Hindi for to put in, then the Hindi for cumin, then say \emph{baghārnā} again
\end{itemize}

\begin{tcolorbox}[breakable,colback=teal!4,colframe=teal!35!black,title={Before you move on}]
What does \hi{बघारना} mean? (\textquotedblleft{}To temper (with hot spiced oil)\textquotedblright{}.) Say it once more.
\end{tcolorbox}
\section[\texorpdfstring{\hi{भूँजना} (bhū̃jnā)}{भूँजना (bhū̃jnā)}]{\hi{भूँजना} (bhū̃jnā) — to roast; to parch (dry, in a pan)}

\label{lesson:HI-C521-bhujna}



\begin{quote}
Before the new one: say the Hindi for cumin, then the Hindi for to temper (with hot spiced oil).
\end{quote}

\subsection*{You'll want to know: \hi{भूँजना}}
\textbf{\hi{भूँजना}} — \emph{bhū̃jnā} — \textquotedblleft{}to roast; to parch (dry, in a pan)\textquotedblright{}.

\begin{quote}
\hi{दादी} \hi{कड़ाही} \hi{में} \hi{मूँगफली} \hi{भूँज} \hi{रही} \hi{हैं।}
\end{quote}

\begin{quote}
\emph{Dādī kaṛāhī meṁ mū̃gphalī bhū̃j rahī haiṁ.}
\end{quote}

\begin{quote}
Grandmother is roasting peanuts in the wok.
\end{quote}

Past, with \hi{ने} (he or she …): \hi{उसने} \hi{मूँगफली} \hi{भूँजी}, \emph{usne mū̃gphalī bhū̃jī}.

\begin{cousinweb}
\textbf{\hi{भूँजना}} (\emph{bhū̃jnā}) is a tadbhava, inherited from \emph{*bhṛñjati}, \textquotedblleft{}parches\textquotedblright{}, a form CDIAL reconstructs from its descendants (CDIAL 9586). CDIAL gives the Hindi verb as \textquotedblleft{}to parch, fry\textquotedblright{}.
\end{cousinweb}

\begin{grammarlens}[title={What you've built}]
One more word for describing what you lived and felt.
\end{grammarlens}

\subsection*{Your turn}
\begin{itemize}
\raggedright
  \item \emph{Say it:} \emph{bhū̃jnā}
  \item \emph{Say it:} \emph{bhū̃jnā}, once more
  \item \emph{Say it:} \emph{Dādī kaṛāhī meṁ mū̃gphalī bhū̃j rahī haiṁ.} — Grandmother is roasting peanuts in the wok.
  \item \emph{Recall:} say the Hindi for cumin, then the Hindi for to temper (with hot spiced oil), then say \emph{bhū̃jnā} again
\end{itemize}

\begin{tcolorbox}[breakable,colback=teal!4,colframe=teal!35!black,title={Before you move on}]
What does \hi{भूँजना} mean? (\textquotedblleft{}To roast; to parch (dry, in a pan)\textquotedblright{}.) Say it once more.
\end{tcolorbox}
\section[(dialogue)]{Looking back: from Metal and Gold to Heat in the Kitchen}

\label{lesson:HI-R521-review-514-521}



\begin{quote}
Say the Hindi for a stamp, a die and to file (metal or wood).
\end{quote}

\subsection*{Your turn}
Say these aloud:

\begin{itemize}
\raggedright
  \item a stamp, a die, to file (metal or wood), to set (a stone), to polish, to make shine, to melt (something)
  \item loose, a screw, a screwdriver, to tighten, to rattle
  \item to let drip, to crack, to split, a pipe, a hole, to trickle down
  \item a brick, lime (for whitewash), plaster (on a wall), to smear (a floor or wall with mud), to lay (bricks)
  \item a wire, a light bulb, a switch, a peg, to hang up
  \item a trunk, a spider, waste paper, to sweep (with a broom), to scour (pots)
  \item to kindle, to knead, a round board for rolling dough, to roll out (dough), to puff up
  \item the heat of a flame, to put in, cumin, to temper (with hot spiced oil), to roast
\end{itemize}

\begin{tcolorbox}[breakable,colback=teal!4,colframe=teal!35!black,title={Before you move on}]
A stamp, a die? (\textbf{\hi{ठप्पा}}, \emph{ṭhappā}.) To file (metal or wood)? (\textbf{\hi{रेतना}}, \emph{retnā}.) To set (a stone)? (\textbf{\hi{जड़ना}}, \emph{jaṛnā}.) To polish, to make shine? (\textbf{\hi{चमकाना}}, \emph{camkānā}.) To melt (something)? (\textbf{\hi{पिघलाना}}, \emph{pighlānā}.) Loose? (\textbf{\hi{ढीला}}, \emph{ḍhīlā}.) A screw? (\textbf{\hi{पेंच}}, \emph{peñc}.) A screwdriver? (\textbf{\hi{पेचकस}}, \emph{peckas}.) To tighten? (\textbf{\hi{कसना}}, \emph{kasnā}.) To rattle? (\textbf{\hi{खड़खड़ाना}}, \emph{khaṛkhaṛānā}.) To let drip? (\textbf{\hi{टपकाना}}, \emph{ṭapkānā}.) To crack, to split? (\textbf{\hi{दरकना}}, \emph{daraknā}.) A pipe? (\textbf{\hi{पाइप}}, \emph{pāip}.) A hole? (\textbf{\hi{छेद}}, \emph{ched}.) To trickle down? (\textbf{\hi{ढलकना}}, \emph{ḍhalaknā}.) A brick? (\textbf{\hi{ईंट}}, \emph{ī̃ṭ}.) Lime (for whitewash)? (\textbf{\hi{चूना}}, \emph{cūnā}.) Plaster (on a wall)? (\textbf{\hi{पलस्तर}}, \emph{palastar}.) To smear (a floor or wall with mud)? (\textbf{\hi{लीपना}}, \emph{līpnā}.) To lay (bricks)? (\textbf{\hi{चिनना}}, \emph{cinnā}.) A wire? (\textbf{\hi{तार}}, \emph{tār}.) A light bulb? (\textbf{\hi{बल्ब}}, \emph{balb}.) A switch? (\textbf{\hi{स्विच}}, \emph{svic}.) A peg? (\textbf{\hi{खूँटी}}, \emph{khū̃ṭī}.) To hang up? (\textbf{\hi{टाँगना}}, \emph{ṭā̃gnā}.) A trunk? (\textbf{\hi{संदूक़}}, \emph{sandūq}.) A spider? (\textbf{\hi{मकड़ी}}, \emph{makṛī}.) Waste paper? (\textbf{\hi{रद्दी}}, \emph{raddī}.) To sweep (with a broom)? (\textbf{\hi{बुहारना}}, \emph{buhārnā}.) To scour (pots)? (\textbf{\hi{माँजना}}, \emph{mā̃jnā}.) To kindle? (\textbf{\hi{सुलगाना}}, \emph{sulgānā}.) To knead? (\textbf{\hi{गूँधना}}, \emph{gū̃dhnā}.) A round board for rolling dough? (\textbf{\hi{चकला}}, \emph{caklā}.) To roll out (dough)? (\textbf{\hi{बेलना}}, \emph{belnā}.) To puff up? (\textbf{\hi{फूलना}}, \emph{phūlnā}.) The heat of a flame? (\textbf{\hi{आँच}}, \emph{ā̃c}.) To put in? (\textbf{\hi{डालना}}, \emph{ḍālnā}.) Cumin? (\textbf{\hi{जीरा}}, \emph{jīrā}.) To temper (with hot spiced oil)? (\textbf{\hi{बघारना}}, \emph{baghārnā}.) To roast? (\textbf{\hi{भूँजना}}, \emph{bhū̃jnā}.)
\end{tcolorbox}
